% GENERATED FILE. Edit canonical lessons, then run npm run generate:books.
% canonical-source-hash: fnv1a64:b02232c7db7f1acd
% canonical-lessons: SA-C547-utthanam, SA-C547-avanatih, SA-C547-patanam, SA-C547-arohah, SA-C547-ucchalati

\chapter{Ups and Downs}
\label{ch:sa-ups-and-downs}
\hlchaptermodality{547}

\begin{chapteropening}
\textbf{By the end of this chapter:} \emph{I can talk about rising, falling and climbing.}

Complete the last lesson of chapter 547: I can talk about rising, falling and climbing.
\end{chapteropening}

\section[\texorpdfstring{\sk{उत्थानम्} (utthānam)}{उत्थानम् (utthānam)}]{\sk{उत्थानम्} (utthānam) — a rise}

\label{lesson:SA-C547-utthanam}



\begin{quote}
Before the new one: say the Sanskrit for an attainment, then the Sanskrit for a consequence.
\end{quote}

\subsection*{You'll want to know: \sk{उत्थानम्}}
\textbf{\sk{उत्थानम्}} — \emph{utthānam} — \textquotedblleft{}a rise\textquotedblright{}.

\begin{quote}
\sk{यदि} \sk{देशस्य} \sk{उत्थानं} \sk{भवति}, \sk{तर्हि} \sk{जनाः} \sk{प्रसन्नाः} \sk{भविष्यन्ति।}
\end{quote}

\begin{quote}
\emph{yadi deśasya utthānaṃ bhavati, tarhi janāḥ prasannāḥ bhaviṣyanti.}
\end{quote}

\begin{quote}
If the country rises, the people will be glad.
\end{quote}

A neuter noun.

\begin{cousinweb}
\textbf{\sk{उत्थानम्}} (\emph{utthānam}) is built on \emph{ud}, \textquotedblleft{}up\textquotedblright{}, and the root √sthā, \textquotedblleft{}to stand\textquotedblright{}, which also gives \emph{uttiṣṭhati}; MW gives \emph{utthāna} as standing up, rising (MW).
\end{cousinweb}

\begin{grammarlens}[title={What you've built}]
One more word for saying what happens if.
\end{grammarlens}

\subsection*{Your turn}
\begin{itemize}
\raggedright
  \item \emph{Say it:} \emph{utthānam}
  \item \emph{Say it:} \emph{utthānam}, once more
  \item \emph{Say it:} \emph{yadi deśasya utthānaṃ bhavati, tarhi janāḥ prasannāḥ bhaviṣyanti.} — If the country rises, the people will be glad.
  \item \emph{Recall:} say the Sanskrit for an attainment, then the Sanskrit for a consequence, then say \emph{utthānam} again
\end{itemize}

\begin{tcolorbox}[breakable,colback=teal!4,colframe=teal!35!black,title={Before you move on}]
What does \sk{उत्थानम्} mean? (\textquotedblleft{}A rise\textquotedblright{}.) Say it once more.
\end{tcolorbox}
\section[\texorpdfstring{\sk{अवनतिः} (avanatiḥ)}{अवनतिः (avanatiḥ)}]{\sk{अवनतिः} (avanatiḥ) — a decline}

\label{lesson:SA-C547-avanatih}



\begin{quote}
Before the new one: say the Sanskrit for a consequence, then the Sanskrit for a rise.
\end{quote}

\subsection*{You'll want to know: \sk{अवनतिः}}
\textbf{\sk{अवनतिः}} — \emph{avanatiḥ} — \textquotedblleft{}a decline\textquotedblright{}.

\begin{quote}
\sk{यदि} \sk{व्यापारस्य} \sk{अवनतिः} \sk{भवति}, \sk{तर्हि} \sk{वणिक्} \sk{दुःखी} \sk{भवति।}
\end{quote}

\begin{quote}
\emph{yadi vyāpārasya avanatiḥ bhavati, tarhi vaṇik duḥkhī bhavati.}
\end{quote}

\begin{quote}
If trade declines, the merchant is unhappy.
\end{quote}

A feminine noun.

\begin{cousinweb}
\textbf{\sk{अवनतिः}} (\emph{avanatiḥ}) is built on \emph{ava}, \textquotedblleft{}down\textquotedblright{}, and the root √nam, \textquotedblleft{}to bow\textquotedblright{}, which also gives \emph{namati}; MW gives \emph{avanati} as bowing down, stooping; \textquotedblleft{}a decline\textquotedblright{}, of trade or fortune, is the modern sense (MW).
\end{cousinweb}

\begin{grammarlens}[title={What you've built}]
One more word for saying what happens if.
\end{grammarlens}

\subsection*{Your turn}
\begin{itemize}
\raggedright
  \item \emph{Say it:} \emph{avanatiḥ}
  \item \emph{Say it:} \emph{avanatiḥ}, once more
  \item \emph{Say it:} \emph{yadi vyāpārasya avanatiḥ bhavati, tarhi vaṇik duḥkhī bhavati.} — If trade declines, the merchant is unhappy.
  \item \emph{Recall:} say the Sanskrit for a consequence, then the Sanskrit for a rise, then say \emph{avanatiḥ} again
\end{itemize}

\begin{tcolorbox}[breakable,colback=teal!4,colframe=teal!35!black,title={Before you move on}]
What does \sk{अवनतिः} mean? (\textquotedblleft{}A decline\textquotedblright{}.) Say it once more.
\end{tcolorbox}
\section[\texorpdfstring{\sk{पतनम्} (patanam)}{पतनम् (patanam)}]{\sk{पतनम्} (patanam) — a fall}

\label{lesson:SA-C547-patanam}



\begin{quote}
Before the new one: say the Sanskrit for a rise, then the Sanskrit for a decline.
\end{quote}

\subsection*{You'll want to know: \sk{पतनम्}}
\textbf{\sk{पतनम्}} — \emph{patanam} — \textquotedblleft{}a fall\textquotedblright{}.

\begin{quote}
\sk{यदि} \sk{पतनं} \sk{भवति}, \sk{तर्हि} \sk{पुनः} \sk{उत्तिष्ठतु।}
\end{quote}

\begin{quote}
\emph{yadi patanaṃ bhavati, tarhi punaḥ uttiṣṭhatu.}
\end{quote}

\begin{quote}
If you fall, get up again.
\end{quote}

A neuter noun.

\begin{cousinweb}
\textbf{\sk{पतनम्}} (\emph{patanam}) comes from the root √pat, \textquotedblleft{}to fall\textquotedblright{}, which also gives \emph{patati}, with the ending \emph{-ana}; MW gives \emph{patana} as flying, falling (MW).
\end{cousinweb}

\begin{grammarlens}[title={What you've built}]
One more word for saying what happens if.
\end{grammarlens}

\subsection*{Your turn}
\begin{itemize}
\raggedright
  \item \emph{Say it:} \emph{patanam}
  \item \emph{Say it:} \emph{patanam}, once more
  \item \emph{Say it:} \emph{yadi patanaṃ bhavati, tarhi punaḥ uttiṣṭhatu.} — If you fall, get up again.
  \item \emph{Recall:} say the Sanskrit for a rise, then the Sanskrit for a decline, then say \emph{patanam} again
\end{itemize}

\begin{tcolorbox}[breakable,colback=teal!4,colframe=teal!35!black,title={Before you move on}]
What does \sk{पतनम्} mean? (\textquotedblleft{}A fall\textquotedblright{}.) Say it once more.
\end{tcolorbox}
\section[\texorpdfstring{\sk{आरोहः} (ārohaḥ)}{आरोहः (ārohaḥ)}]{\sk{आरोहः} (ārohaḥ) — a climb; an ascent}

\label{lesson:SA-C547-arohah}



\begin{quote}
Before the new one: say the Sanskrit for a decline, then the Sanskrit for a fall.
\end{quote}

\subsection*{You'll want to know: \sk{आरोहः}}
\textbf{\sk{आरोहः}} — \emph{ārohaḥ} — \textquotedblleft{}a climb; an ascent\textquotedblright{}.

\begin{quote}
\sk{यदि} \sk{पर्वतस्य} \sk{आरोहः} \sk{कठिनः} \sk{अस्ति}, \sk{तर्हि} \sk{मन्दं} \sk{गच्छतु।}
\end{quote}

\begin{quote}
\emph{yadi parvatasya ārohaḥ kaṭhinaḥ asti, tarhi mandaṃ gacchatu.}
\end{quote}

\begin{quote}
If the climb up the mountain is hard, go slowly.
\end{quote}

A masculine noun.

\begin{cousinweb}
\textbf{\sk{आरोहः}} (\emph{ārohaḥ}) is built on \emph{ā} and the root √ruh, \textquotedblleft{}to rise, to grow\textquotedblright{}, which also gives \emph{ārohati}; MW gives \emph{āroha} as one who mounts, and as rising, an ascent (MW).
\end{cousinweb}

\begin{grammarlens}[title={What you've built}]
One more word for saying what happens if.
\end{grammarlens}

\subsection*{Your turn}
\begin{itemize}
\raggedright
  \item \emph{Say it:} \emph{ārohaḥ}
  \item \emph{Say it:} \emph{ārohaḥ}, once more
  \item \emph{Say it:} \emph{yadi parvatasya ārohaḥ kaṭhinaḥ asti, tarhi mandaṃ gacchatu.} — If the climb up the mountain is hard, go slowly.
  \item \emph{Recall:} say the Sanskrit for a decline, then the Sanskrit for a fall, then say \emph{ārohaḥ} again
\end{itemize}

\begin{tcolorbox}[breakable,colback=teal!4,colframe=teal!35!black,title={Before you move on}]
What does \sk{आरोहः} mean? (\textquotedblleft{}A climb; an ascent\textquotedblright{}.) Say it once more.
\end{tcolorbox}
\section[\texorpdfstring{\sk{उच्छलति} (ucchalati)}{उच्छलति (ucchalati)}]{\sk{उच्छलति} (ucchalati) — springs up}

\label{lesson:SA-C547-ucchalati}



\begin{quote}
Before the new one: say the Sanskrit for a fall, then the Sanskrit for a climb.
\end{quote}

\subsection*{You'll want to know: \sk{उच्छलति}}
\textbf{\sk{उच्छलति}} — \emph{ucchalati} — \textquotedblleft{}springs up\textquotedblright{}.

\begin{quote}
\sk{यदि} \sk{जले} \sk{शिला} \sk{पतति}, \sk{तर्हि} \sk{जलम्} \sk{उच्छलति।}
\end{quote}

\begin{quote}
\emph{yadi jale śilā patati, tarhi jalam ucchalati.}
\end{quote}

\begin{quote}
If a stone falls into the water, the water splashes up.
\end{quote}

Past (it …): \sk{जलम्} \sk{उच्छलितवत्}, \emph{jalam ucchalitavat}.

\begin{cousinweb}
\textbf{\sk{उच्छलति}} (\emph{ucchalati}) is built on \emph{ud}, \textquotedblleft{}up\textquotedblright{}, and the root √śal, \textquotedblleft{}to move\textquotedblright{}, which after \emph{ud} is written \emph{chal}; MW gives \emph{ucchal} as to fly upwards, to spring up (MW).
\end{cousinweb}

\begin{grammarlens}[title={What you've built}]
One more word for saying what happens if.
\end{grammarlens}

\subsection*{Your turn}
\begin{itemize}
\raggedright
  \item \emph{Say it:} \emph{ucchalati}
  \item \emph{Say it:} \emph{ucchalati}, once more
  \item \emph{Say it:} \emph{yadi jale śilā patati, tarhi jalam ucchalati.} — If a stone falls into the water, the water splashes up.
  \item \emph{Recall:} say the Sanskrit for a fall, then the Sanskrit for a climb, then say \emph{ucchalati} again
\end{itemize}

\begin{tcolorbox}[breakable,colback=teal!4,colframe=teal!35!black,title={Before you move on}]
What does \sk{उच्छलति} mean? (\textquotedblleft{}Springs up\textquotedblright{}.) Say it once more.
\end{tcolorbox}
